\section{Operational Threshold Calibration: Translating Scores into Business Decisions}
\label{sec:threshold_calibration}

A pervasive failure mode in industrial computer vision deployments is the uncalibrated transfer of decision thresholds derived during laboratory model development directly onto the factory floor. In academic literature, algorithms are benchmarked using integrated area-under-the-curve metrics such as AUROC, PR-AUC, or AUPIMO. These metrics compress discriminatory capability across the entire operating spectrum into a single scalar. Whilst indispensable for offline algorithm selection -- resolving which mathematical architecture achieves superior feature separation between nominal and defective populations -- these metrics provide zero guidance for operational commissioning.

An industrial quality gate does not execute a continuous characteristic curve; it executes an instantaneous binary divert decision. At any given conveyor cycle, a manufactured component is either passed downstream or diverted into a secondary quarantine buffer. The numerical threshold governing this decision represents the single most sensitive operational parameter in the facility, and it must be calibrated to minimise total economic loss rather than to maximise an academic abstraction.

\subsection{The Operational Distinction Between AUPIMO and Decision Thresholds}

The operational distinction between AUPIMO and line-side decision thresholds warrants plain explanation. Think of AUPIMO as a fixed benchmark of optical vision quality: it measures how cleanly and sharply the system outlines defects under strict low-false-alarm conditions, completely independent of where the pass/divert threshold is set. In short, it tells plant management whether the underlying model has the visual clarity to spot subtle scratches without flagging clean surfaces.

However, AUPIMO itself does not tell a conveyor mechanism when to push a part into a reject bin. That requires an operational decision threshold calibrated to the factory's specific economic balance. A vision engine with outstanding inherent optical sharpness ($\text{AUPIMO} = 0.603$) will fail commercially if coupled to an arbitrary, uncalibrated threshold, whereas a model with an AUPIMO of $0.500$ calibrated precisely to the customer's actual costs of false alarms versus missed defects will deliver immediate operational value.

By the same token, pixel-level AUROC is categorically discarded as an operational metric. Because conforming pixels make up roughly $98.5\%$ of total component surface area, AUROC is overwhelmed by the massive pool of easy normal pixels, producing deceptively high scores regardless of true defect sensitivity. While reported in Table~\ref{tab:paradigm_comparison} to ensure benchmark reproducibility, it exerts zero influence on operational threshold calibration.

\subsection{Strict Zero-Test-Leakage Calibration Protocol}

The platform mandates an immutable, audit-compliant calibration protocol that enforces strict isolation between model development, threshold derivation, and final acceptance testing.

For every component category, all available nominal training imagery is partitioned deterministically (seed 42) into an 85\% fitting partition, dedicated to coreset memory bank construction, and a 15\% normal validation partition, held out exclusively for statistical threshold calibration. Crucially, the official evaluation split (the test set) remains completely untouched and unseen throughout the entire fitting and calibration lifecycle. This strict partition enforces authentic commissioning conditions: in a live factory rollout, engineering teams possess access solely to nominal production samples, with defect examples remaining entirely unobserved prior to production deployment.

The operational decision threshold $\tau$ is derived empirically from the score distribution generated across the held-out normal validation partition. Formally,

\begin{equation}
    \tau_{\alpha} = \mathrm{Quantile}_{1 - \alpha}\!\left(\{s_i \mid x_i \in \mathcal{D}_{\mathrm{val, normal}}\}\right),
    \label{eq:threshold_quantile}
\end{equation}

where $\alpha$ denotes the plant's tolerable false-alarm budget (the Type I error ceiling). Specifying $\alpha = 0.005$ establishes an empirical threshold ensuring that $99.5\%$ of conforming components within the calibration population pass the gate without intervention; conversely, specifying $\alpha = 0.05$ accepts a 5\% pseudo-scrap rate to guarantee aggressive defect capture. As established in Section~\ref{sec:cost_model}, the parameter $\alpha$ is governed strictly by the client's financial cost ratio $\rho$ via Equation~\ref{eq:optimal_threshold}.

\subsection{Standardised Risk-Tiered Operating Profiles}

To operationalise this framework across diverse manufacturing sectors, the platform standardises three predefined risk tiers:

\begin{itemize}[leftmargin=1.5em]
    \item \textbf{Tier 1: Life-Safety and High-Liability (e.g., Pharmaceuticals, Medical Devices, Critical Semiconductors):} Governed by an $F_2$-optimised calibration objective ($\alpha \approx 0.05$). This profile weights defect recall four times more heavily than precision, shifting the decision threshold into the nominal distribution tail. Every borderline component is diverted to manual audit, eliminating defect escapes at the expense of a controlled, quantified pseudo-scrap rate.
    \item \textbf{Tier 2: Precision Industrial (e.g., Automotive Powertrain, Bearing Assemblies, Structural Cables):} Governed by a balanced $F_1$-optimised objective ($\alpha \approx 0.01$). This profile establishes a cost-neutral equilibrium between Type I and Type II errors, appropriate for mid-tier risk environments where defect escapes incur rework rather than total product recalls.
    \item \textbf{Tier 3: High-Throughput Commodity (e.g., Threaded Fasteners, Stamped Washers, Continuous Zippers):} Governed by an $F_{0.5}$-optimised objective ($\alpha \approx 0.002$). This profile penalises false alarms twice as heavily as defect escapes, preventing pseudo-scrap from choking downstream automated packaging equipment while reliably catching severe structural anomalies.
\end{itemize}

Figure~\ref{fig:threshold_calibration} depicts the empirical anomaly score distributions across normal validation and defective evaluation sets, illustrating how the physical separation between distributions dictates achievable defect containment across the three risk tiers.

\begin{figure}[tp]
    \centering
    \placeholder{2.2cm}{Two-panel or multi-panel figure showing empirical anomaly score distributions. Left panel (or top row): kernel density estimates (KDE) or histograms of image-level anomaly scores for normal validation images (shaded in blue/grey) and anomalous test images (shaded in red/orange) for a low-difficulty category (e.g., Leather or Bottle). Right panel (or bottom row): same for a high-difficulty category (e.g., Screw or Cable), where the distributions substantially overlap. Vertical dashed lines mark the three operational threshold tiers ($\tau_{0.002}$, $\tau_{0.01}$, $\tau_{0.05}$). The figure should make visually clear that threshold choice directly controls the precision-recall trade-off and that the achievable recall depends on the degree of score distribution separation.}
    \caption{Empirical anomaly score distributions on the normal validation partition (blue) and the anomalous evaluation split (red) for representative categories. Vertical dashed lines mark the three operational threshold tiers calibrated under our strict zero-test-leakage protocol. Categories exhibiting sharp distribution separation (e.g., Leather) permit aggressive thresholds that attain zero escape with negligible false alarms; categories with overlapping distributions (e.g., Screw) necessitate explicit economic optimisation to establish the operational operating point.}
    \label{fig:threshold_calibration}
\end{figure}
